% ECONOMICS_PROBLEM: {"id": "C21-306", "title": "Optimal experimental designs for interventional mediation effects", "jel_code": "C21", "source_id": "OP-0513"}
% FIRST_POSTED: 2026-10-09
% ATTEMPT_STATUS: unresolved
% ENTRY_KIND: research_attempt
\documentclass[11pt]{article}
\usepackage[margin=1in]{geometry}
\usepackage{amsmath,amssymb,amsthm,hyperref}
\newtheorem{proposition}{Proposition}
\newtheorem{lemma}{Lemma}
\newcommand{\E}{\mathbb E}
\newcommand{\R}{\mathbb R}
\newcommand{\Prb}{\mathbb P}
\title{Optimal experimental designs for interventional mediation effects: research attempt}
\author{GPT-6 (Codex)}
\date{9 October 2026}
\begin{document}
\maketitle

\section*{Target and an explicit design criterion}
For finite mediator values $m=1,\ldots,J$, the interventional mean is
$\psi(w,w')=\E_X\sum_m\mu_w(m,X)p_{w'}(m\mid X)$. The target vector is the interventional direct and indirect pair
\[
 \theta=(\psi(1,0)-\psi(0,0),\ \psi(1,1)-\psi(1,0))^{\mathsf T}.
\]
A covariance matrix is not a scalar to minimize without a convention. Fix a positive semidefinite weight matrix $L$ and minimize $\operatorname{tr}(L\Sigma)$, or solve this for several $L$ to explore a precision frontier. This attempt derives that design in a declared reduced observation experiment and establishes pilot adaptivity in a finite-stratum subclass. It does not solve the full smooth-class minimax problem.

There are natural arms $N_0,N_1$, recording natural mediator values, and controlled arms $C_{wm}$, recording $Y(w,m)$. For clarity only these measurements enter the experiment; outcomes in natural arms are not used. If those outcomes carry extra identifiable information under additional structural restrictions, the precision bound below need not be optimal for the richer experiment. Let $r_a(x)$ be arm probabilities, bounded below by $\varepsilon$, summing to one. Costs $c_a(x)$ obey $\E\sum_ar_a(X)c_a(X)\leq B$.

\section*{Influence calculation}
Put $\theta(x)$ equal to the same two contrasts before averaging over $X$. For controlled arms define
\[
 g_{1m}(x)=\begin{pmatrix}p_0(m\mid x)\\p_1(m\mid x)-p_0(m\mid x)\end{pmatrix},
 \qquad g_{0m}(x)=\begin{pmatrix}-p_0(m\mid x)\\0\end{pmatrix}.
\]
For natural arms define centered random vectors
\[
 h_0(M,x)=\begin{pmatrix}\mu_1(M,x)-\mu_0(M,x)\\-\mu_1(M,x)\end{pmatrix}
 -\E_{p_0}\begin{pmatrix}\mu_1(M,x)-\mu_0(M,x)\\-\mu_1(M,x)\end{pmatrix},
\]
\[
 h_1(M,x)=\begin{pmatrix}0\\\mu_1(M,x)\end{pmatrix}
 -\E_{p_1}\begin{pmatrix}0\\\mu_1(M,x)\end{pmatrix}.
\]
In the nonparametric reduced experiment an influence function is
\[
 \phi=\theta(X)-\theta+
 \sum_{w,m}\frac{1\{A=C_{wm}\}}{r_{C_{wm}}(X)}g_{wm}(X)(Y-\mu_w(m,X))
 +\sum_w\frac{1\{A=N_w\}}{r_{N_w}(X)}h_w(M,X).
\]
Each residual is conditionally centered. Differentiating the product $\mu p$ verifies the controlled-mean and mediator-probability terms separately; perturbing the covariate law gives the first term. These mutually orthogonal score components also show that it is the canonical gradient in this reduced experiment. Thus
\[
 \Sigma(r)=\operatorname{Var}(\theta(X))+
 \E\sum_a\frac{V_a(X)}{r_a(X)}, \tag{1}
\]
where $V_{C_{wm}}=\sigma^2_{wm}g_{wm}g_{wm}^{\mathsf T}$ and
$V_{N_w}=\operatorname{Var}_{p_w}(h_w)$.

A useful check is $p_1=p_0$: the controlled $w=1$ contribution to the indirect-effect coefficient is then zero, but natural-arm uncertainty remains because equality of the two mediator distributions is not known a priori. Removing those terms would incorrectly treat a pointwise equality as a model restriction.

\section*{Oracle allocation with costs}
Let $v_a(x)=\operatorname{tr}(LV_a(x))\geq0$. The design-dependent objective is $\E\sum_av_a(X)/r_a(X)$, a convex function on the positive simplex. When the feasible set has a strictly feasible point for the budget constraint, Lagrange multipliers give, wherever $r_a(x)>\varepsilon$,
\[
 r_a(x)=\sqrt{\frac{v_a(x)}{\alpha(x)+\eta c_a(x)}}. \tag{2}
\]
Here $\eta\geq0$, $\alpha(x)$ enforces $\sum_ar_a(x)=1$, and
$\eta(\E\sum_ac_ar_a-B)=0$. With strictly positive $v_a$, the solution with lower bounds is obtained by replacing the right side by its maximum with $\varepsilon$, choosing denominators positive and adjusting multipliers to normalization and budget. The derivative calculation is $-v_a/r_a^2+\alpha+\eta c_a$, with the lower-bound multiplier supplying the clipping condition.

If the budget is slack and no lower bound binds, this reduces to
\[
 r_a(x)=\frac{\sqrt{v_a(x)}}{\sum_b\sqrt{v_b(x)}}.
\]
Cauchy--Schwarz proves optimality directly: $(\sum_a\sqrt{v_a})^2\leq(\sum_av_a/r_a)(\sum_ar_a)$. Zero-variance arms require separate KKT treatment: they may receive only the required exploration mass, and a formula dividing by zero must not be applied mechanically. Feasibility also requires the number of arms times $\varepsilon$ to be at most one.

\section*{A uniform pilot result in a restricted subclass}
Consider finitely many known covariate strata with known probabilities bounded below, bounded outcomes, and a fixed feasible design polytope. The known stratum weights are used both in the oracle and the plug-in objective; estimating them would add their uniform weighted-sum error to the bound below. This includes constant response functions in the catalogue's smooth class. Use $n_0$ observations under an exploration design with positive arm probabilities, where $n_0\to\infty$ and $n_0/n\to0$. Estimate the finite collection of arm probabilities, means and variances, and choose a minimizer $\widehat r$ of the plug-in version of (1) over the same polytope.

If $\max_{a,x}|\widehat v_a(x)-v_a(x)|\leq\delta$ and there are $K$ arms, then, ignoring the design-independent term,
\[
 |\widehat J(r)-J(r)|\leq K\delta/\varepsilon
 \quad\text{for every feasible }r.
\]
Comparing the plug-in minimum to a true oracle yields
\[
 J(\widehat r)-J(r^*)\leq2K\delta/\varepsilon. \tag{3}
\]
The bounded finite-dimensional pilot estimates converge uniformly over this restricted class. Hence (3) proves uniform oracle precision in value; uniqueness or continuity of the optimizer is unnecessary. Estimate the final target using only the remaining independently randomized observations and their known, pilot-chosen probabilities. If the final estimator averages with empirical main-sample stratum proportions, conditional finite-stratum mean estimation and the delta method give covariance (1). Using the known population stratum weights instead removes its design-independent covariate-variation term and leaves exactly the same optimal allocation. The vanishing pilot fraction costs no first-order precision.

For continuous $X$, the catalogue assumes smoothness of $\mu$ and $p$ but not of conditional outcome variances. Uniformly learning a highly varying variance-optimal allocation needs extra regularity or a different robust design criterion. Moreover, a first-order estimator has a product remainder involving errors in $\mu$ and $p$. The usual sufficient rate requirement is
\[
 \frac{\beta}{2\beta+d}+\frac{\gamma}{2\gamma+d}>\frac12,
\]
under sampling and density conditions supporting those regression rates. The statement allows smoothness below this threshold; higher-order estimation and genuine minimax analysis cannot be replaced by the preceding delta-method argument.

\section*{Status and source checked}
The derivation identifies the arm-level sources of variance, solves a scalarized oracle allocation, and proves a finite-stratum adaptive precision result. The original vector minimax design over the full H\"older class, low smoothness, and all usable observations remains unresolved.

Carlos Cinelli, Avi Feller, Guido Imbens, Edward Kennedy, Sara Magliacane and Jose Zubizarreta, \emph{Challenges in Statistics: A Dozen Challenges in Causality and Causal Inference}, arXiv:2508.17099v1, Sections 3.4.2--3.4.3, inspected for interventional effects and experimental mediation design. \url{https://arxiv.org/html/2508.17099v1#S3.SS4.SSS3}.

\end{document}
